\begin{document}
\maketitle

\begin{abstract}
This paper isolates the geometric and metrological branch of TFPT. The theory is not presented as
predicting isolated SI numbers. Instead it constructs an internal dimensionless metrology from the
boundary spectral unit $\lambda_\Sigma$, with gravity, Planck normalization, electroweak matching,
and pole readouts expressed as boundary-normalized observables.
\end{abstract}

\begin{tfptfrontbox}
\textbf{Inputs from previous papers.} Paper 1 supplies the boundary branch and primitive spectral
unit. Paper 2 supplies the carrier/Higgs structure. Paper 4 may supply the renormalized observable
layer when the analytic QFT closure is referenced.

\textbf{New theorem contribution.} The boundary-normalized metrology:
\[
\lambda_\Sigma=\lambda_1^+(|B_\Sigma|),\qquad
\rho_\star=\frac{\chigeozero^2}{\lambda_\Sigma^2},\qquad
\frac{\Mpl^2}{\lambda_\Sigma^2}=\frac{\rho_\star}{2\pi^2},\qquad
G_N\lambda_\Sigma^2=\frac{\pi}{4\rho_\star},
\]
and the Einstein-limit normalizer
\[
\xi=\frac{\cthree}{\phiz},
\qquad
\kappa^2=\xi\,\frac{\phiz}{\cthree^2},
\qquad
\xi_{\mathrm{tree}}=\tfrac{3}{4},\;\;
\xi_\star\approx 0.748327808\ldots
\]

\textbf{Not claimed here.} No late-time $H_0$, no CMB, no black holes or horizons, no E8 stage
atlas, and no astrophysical bursts.

\textbf{Falsification or audit surface.} The paper fails if SI units enter as hidden inputs, if
$\lambda_\Sigma$ is not fixed by the boundary branch, or if electroweak matching is mixed with
cosmological comparison rows.
\end{tfptfrontbox}


\begin{tfptclaimcontract}
\textbf{Claim.} Boundary spectral unit $\lambda_\Sigma$ defines dimensionless metrology for gravity and EW readouts.\par
\textbf{Inputs.} Boundary spectral unit, geometric Hodge branch, renormalized observable layer.\par
\textbf{First assumptions.} Stationary root for $\rho_\star$, boundary-normalized functional, scheme projection only after physical readout.\par
\textbf{Proof status.} Bridge/metrology theorem target.\par
\textbf{Kill condition.} Any hidden SI input such as external $G_N$, $G_F$, GeV scale, or Planck normalization enters before the readout map.\par
\end{tfptclaimcontract}

\begin{tfptstatusbox}
\textbf{\stag{L} (Lean verified) used as input.}\ The $3+2$
hypercharge carrier of Paper 2 enters as the algebraic anchor
for the boundary-normalised metrology.

\smallskip
\textbf{\stag{I} (Lean interfaced).}\ The Paper-1 input enters
as \texttt{BoundaryPolarization}. The Higgs rank-2 interface
enters via \texttt{HiggsTopForm} (Lean-proved at the algebraic
shadow level: $E_+ \simeq K[X_0,X_1]_1$).

\smallskip
\textbf{Higgs-rank-2 origin: \stag{M} with \stag{L} algebraic
shadow and \stag{S} Proj cohomology.}\ The Higgs rank-2
discharge is reached by combining three pieces:
\begin{itemize}[leftmargin=1.6em, itemsep=0pt]
    \item the Lean theorem
        \texttt{HiggsSchemeCohomologyShadow.toDegreeOneTwoVarForm}
        ($K[X_0,X_1]_1 \simeq_K \texttt{DegreeOneTwoVarForm}\,K$,
        \stag{L});
    \item the present paper's manuscript proof of the compact
        Higgs index on the seam-even bosonic block (\stag{M});
    \item the standard algebraic-geometry identification
        $H^0(\mathrm{Proj}\,K[X,Y], \mathcal O(1)) \simeq K[X,Y]_1$
        (\stag{S}: cited, not reproved).
\end{itemize}
Taking those three together, the Higgs rank-2 statement
$\dim E_+ = 2$ is therefore \stag{M} as a TFPT theorem with a
Lean-verified algebraic shadow.

\smallskip
\textbf{\stag{C} (downstream observable readouts).}\ The
boundary spectral unit $\lambda_\Sigma$ as a Planck
normalisation, the resulting electroweak matching, and the
dimensionless observable functor are presented as
\emph{conditional} downstream readouts. The audit obligation
is that no SI input enters before the readout map.
\end{tfptstatusbox}

\tableofcontents

\section{Boundary Spectral Unit}

The internal unit is
\[
\lambda_\Sigma=\lambda_1^+(|B_\Sigma|).
\]
All dimensionful-looking statements in this paper are rewritten as dimensionless quotients by
powers of $\lambda_\Sigma$. This is the editorial rule that separates metrology from numerology.

\section{Geometric Hodge Closure}

The geometric branch reconstructs a spectral Einstein functional and its Hodge closure. The paper
should retain the source draft's geometric reconstruction, spectral Einstein functional, and Planck
closure, but it should remove cosmological and horizon continuations.

\section{Planck Normalization}

The Planck readout is internal:
\[
\rho_\star=\frac{\chigeozero^2}{\lambda_\Sigma^2},
\qquad
\frac{\Mpl^2}{\lambda_\Sigma^2}=\frac{\rho_\star}{2\pi^2},
\qquad
G_N\lambda_\Sigma^2=\frac{\pi}{4\rho_\star}.
\]
The question is not which SI value of $G_N$ is inserted, but which dimensionless branch quotient is
fixed.

\subsection{Einstein-limit normalizer \texorpdfstring{$\xi=\cthree/\phiz$}{xi = c3/phi0}}
\label{ssec:einstein-limit-normalizer}

Once the boundary spectral unit $\lambda_\Sigma$, the seam constant $\cthree$, and the retained
seed $\phiz$ are fixed, a single dimensionless quotient governs the transition between the
UFE-normalized boundary functional and the Einstein--Hilbert presentation. Writing
$\kappa^2=8\pi G_N$ for the Einstein--Hilbert prefactor, the boundary-branch reduction yields a
relation of the form
\begin{equation}
\label{eq:kappa-xi}
\kappa^2 \;=\; \xi\,\frac{\phiz}{\cthree^2},
\qquad
\xi \;=\; \frac{\cthree}{\phiz}.
\end{equation}
The product is intrinsic: only one number, $\xi$, parametrises how the UFE normalization is
rephrased in the Einstein--Hilbert form. At tree level on the canonical branch one finds
\[
\xi_{\mathrm{tree}}=\frac{3}{4},
\]
and the self-consistent geometric reduction gives the closed-branch readout
\[
\xi_\star \;\approx\; 0.748327808\ldots
\]
The two values agree to better than $0.3\%$, which is the structural sense in which the tree
identity $\xi=3/4$ already captures the boundary-normalised Einstein limit.

\begin{tfptfrontbox}
\textbf{What is and is not claimed by $\xi$.} Equation \eqref{eq:kappa-xi} does \emph{not} derive
the SI numerical value of $G_N$ from two dimensionless numbers. It says, in a precise and
boundary-normalised sense, that the dimensionless transition factor between the UFE functional
and the Einstein--Hilbert action is fixed by the same invariant pair $(\cthree,\phiz)$ that fixes
$\alpha$ via the carrier-form closure and $\beta_{\mathrm{rad}}$ via the determinant-line
response. Gravity is not promoted to a primitive observable here. The claim is a compression
identity for the gravitational normaliser inside the dimensionless metrology layer.
\end{tfptfrontbox}

\begin{remark}[Why this is not a fitting knob]
The ratio $\cthree/\phiz$ is not free: $\cthree=1/(8\pi)$ from the boundary primitive (Paper~1)
and $\phiz=\frac{1}{6\pi}+\frac{3}{256\pi^4}$ from the retained seed decoder (Paper~3). Both are
fixed before the gravity readout is touched. Substituting these values into
\eqref{eq:kappa-xi} returns the closed-branch $\xi_\star$ above, which slots into the same
boundary-normalised metrology that already controls $\Mpl^2/\lambda_\Sigma^2$.
\end{remark}

\section{Electroweak Matching}

The electroweak matching layer is included only as boundary-normalized metrology:
\[
v_{\mathrm{phys}}=\vgeo\sqrt{Z_{\mathrm{EW}}^{\mathrm{TFPT}}},
\qquad
G_Nv_{\mathrm{phys}}^2
\]
is an internal dimensionless output. Pole matching may be included in compact form when it is
expressed as a quotient by $\lambda_\Sigma$.

\section{Observable Functor}

The output layer is organized as
\[
\Tstar
\xrightarrow{\mathcal R_{\mathrm{ren}}}
\GrenTFPT
\xrightarrow{\mathcal M_{\mathrm{phys}}}
\OphysTFPT
\xrightarrow{\mathcal M_{\mathrm{scheme}}}
\OschemeTFPT/\SchGrp.
\]
Only the final projection chooses scheme or comparison conventions. The paper should keep this
distinction visible whenever a physical readout is compared to an empirical row.

\clearpage
\section{Main Technical Development}

This section contains the main technical development assigned to this paper by the TFPT 4.5 clean split. Cross-paper background is referenced through dependency and scope boxes; extended backend material is kept in the Technical Companion.

\begin{tfptwarningbox}
\textbf{Status-tier mapping for the source-extract body.}\
The body contains \stag{M} manuscript theorems on the boundary
spectral unit, the absolute spectral Planck closure, the
stationary vacuum theorem, and the Higgs / electroweak matching
on the boundary-normalised functor. None of these are \stag{L}.
Theorems labelled \stag{A} in the legacy source-extract style
should be read as \stag{M}. Downstream observable readouts
(reduced-Planck, Newton, Schwarzschild--de~Sitter, electroweak)
remain \stag{C} until the SI-projection / scheme step is
explicitly separated.
\end{tfptwarningbox}

\tfptDraftXrefNotice

\input{source_extracts/05_metrology_source.tex}

\section{Source Extraction Map}

\begin{tfptsourcebox}
Use \sourceDraft:
\begin{itemize}[leftmargin=1.5em]
    \item Sections 6.4, 6.7, 6.8, and 6.9 for geometric reconstruction, spectral Einstein
    functional, and Planck closure.
    \item Sections 6.21--6.26 for electroweak matching and pole matching.
    \item Section 13.2 only for observable hierarchy parts not already used in Paper 4.
    \item Section 14 in full as the dimensionless metrology core.
    \item Appendix C only in shortened form and without making E8 prominent.
\end{itemize}
\end{tfptsourcebox}

\begin{tfptwarningbox}
Do not include late-time $H_0$, CMB, black holes, horizons, E8 stage atlas, or transient channels.
Those belong to Paper 6 or the Technical Companion.
\end{tfptwarningbox}


\begin{tfptexports}
Exports: $\lambda_\Sigma$, $\rho_\star$, $\bar M_{\mathrm{Pl}}/\lambda_\Sigma$, $G_N\lambda_\Sigma^2$, $v_{\mathrm{phys}}/\lambda_\Sigma$, the Einstein-limit normalizer $\xi=\cthree/\phiz$, and the metrology functor.
\end{tfptexports}

\section{Not Used Here}

Late-time $H_0$, CMB spectra, sky realization, black holes, horizons, E8 stage atlas, transient
channels, and cosmological comparison rows are not used as proof inputs in this paper.

\end{document}
